\documentclass[11pt,a4paper]{article}

\usepackage[margin=2cm]{geometry}
\usepackage[utf8]{inputenc}
\usepackage[T1]{fontenc}
\usepackage{lmodern}
\usepackage{microtype}
\usepackage{hyperref}
\usepackage{booktabs}
\usepackage{enumitem}
\usepackage{xcolor}
\usepackage{titlesec}

\definecolor{accent}{HTML}{E94560}
\hypersetup{colorlinks=true, linkcolor=accent, urlcolor=accent}

\titleformat{\section}{\large\bfseries\sffamily}{}{0pt}{}
\titleformat{\subsection}{\normalsize\bfseries\sffamily}{}{0pt}{}
\titlespacing*{\section}{0pt}{1em}{0.4em}
\titlespacing*{\subsection}{0pt}{0.6em}{0.2em}

\setlength{\parindent}{0pt}
\setlength{\parskip}{0.4em}

\begin{document}

\begin{center}
  {\Large\sffamily\bfseries What 15\,Million Readers Reveal\\About the Human Condition}\\[0.5em]
  {\normalsize COM-480 Data Visualization --- EPFL, Spring 2026 --- Milestone~2}\\[0.3em]
  {\small
    Galanopoulou Rafaila (388285)\\ 
    Migliorelli Claudio (406790)\\
    Poupakis Alexandros (387587)
  }
\end{center}

\vspace{0.2em}
\hrule
\vspace{0.6em}

% ──────────────────────────────────────────────────────────────────
\section{What We Want to Build}

We want to build a scrollytelling website that tells a story about how people read. The idea is simple: take a massive dataset of reading behavior and turn it into something anyone can explore and understand---no statistics background needed.

Our data source is the UCSD Book Graph, which contains 2.36 million books, 229 million user--book interactions, and 15.7 million written reviews. We ask three questions:

\begin{enumerate}[leftmargin=1.5em, itemsep=2pt]
  \item \textbf{Does the world shape what we read?} We look at how genre popularity shifts over time and whether those shifts line up with real-world events like elections or crises.
  \item \textbf{Do readers get pickier over time?} We track how people's average ratings change as they read more books.
  \item \textbf{Do we read what we say we want to read?} We compare how often genres get shelved as ``to-read'' versus how often they actually get read.
\end{enumerate}

The style is inspired by data journalism sites like \textit{The Pudding}---we want the reader to scroll through a narrative where the charts appear naturally as part of the story.

% ──────────────────────────────────────────────────────────────────
\section{Visualizations and Tools}

We have five visualizations planned. The look and feel (colors, fonts, layout) will likely change before the final version.

\begin{table}[h!]
\centering\small
\begin{tabular}{@{}p{0.18\textwidth} p{0.38\textwidth} p{0.18\textwidth} p{0.16\textwidth}@{}}
\toprule
\textbf{Chart} & \textbf{What it shows} & \textbf{Tools} \\
\midrule
Genre Streamgraph &
  How genre shares change quarter by quarter (2010--2017). Hovering highlights one genre. &
  D3.js, Scrollama  \\[3pt]
Event Annotations &
  Vertical markers on the streamgraph for world events (e.g.\ elections, cultural moments). Scroll-triggered. &
  D3.js, Scrollama  \\[3pt]
Rating Drift &
  Average rating vs.\ number of books read, for active users. Shows a downward trend with a confidence band. &
  D3.js \\[3pt]
Aspiration Bars &
  Diverging bar chart: shelved-as-to-read rate on one side, actual reading rate on the other. Sorted by gap. &
  D3.js  \\[3pt]
Summary Counters &
  Animated numbers (books, interactions, reviews, users) that count up as you scroll past them. &
  Vanilla JS  \\
\bottomrule
\end{tabular}
\end{table}

% ──────────────────────────────────────────────────────────────────
\section{What We Need vs.\ What Would Be Nice}

\subsection{Core}
\begin{itemize}[leftmargin=1.5em, itemsep=1pt]
  \item Scrollytelling page with smooth scroll-triggered transitions
  \item Genre streamgraph with event markers and tooltips
  \item Rating drift chart with trend line and confidence band
  \item Aspiration vs.\ reality diverging bars
  \item Animated dataset counters
  \item Data pipeline that turns the raw 11\,GB dataset into small JSON files the website can load
\end{itemize}

\subsection{Nice to have}
\begin{itemize}[leftmargin=1.5em, itemsep=1pt]
  \item Word clouds from review text showing emotional language per genre 
  \item A ``pick a reader type'' explorer showing individual trajectories
  \item Book similarity network clustered by genre 
  \item Play button to animate the streamgraph over time 
  \item Geographic heatmap of genre preferences 
\end{itemize}

% ──────────────────────────────────────────────────────────────────
\section{Current Prototype}

The website is in the \texttt{website/} directory of the repository. It already works with real data---we run a Python script that samples roughly 10\% of the dataset, aggregates it, and writes three JSON files (under 50\,KB total) that D3 loads directly. All five visualizations are functional.

We have deployed our initial website as a GitHub Page in our Github repository. The URL is \url{https://com-480-data-visualization.github.io/CARPA/website/}.

\textbf{What still needs work for Milestone~3:} process the full dataset instead of a sample; write better narrative text for the scroll sections; polish the visual theme; pick and implement one or two of the ``nice to have'' ideas listed above.

% ──────────────────────────────────────────────────────────────────
\section{Handling the Data at Scale}

The full dataset is 11\,GB. The website itself is fast---it only loads tiny pre-aggregated JSON files. The slow part is the one-time preprocessing that creates those files. Right now we sample $\sim$10\% and it takes about 3 minutes. Processing everything with pandas would take 30+ minutes. We are considering three options:

\begin{enumerate}[leftmargin=1.5em, itemsep=2pt]
  \item \textbf{Switch to DuckDB or Polars}---columnar engines that are 5--10$\times$ faster than pandas on this kind of scan-and-aggregate workload.
  \item \textbf{Convert to Parquet first}---a one-time conversion that makes all subsequent queries near-instant.
  \item \textbf{Just sample more}---going from 10\% to 50\% still runs in reasonable time and the aggregated results are statistically stable.
\end{enumerate}

None of these change the website code, only the offline preprocessing step.

\end{document}
